\chapter*{Введение: зачем этот фундамент}
\addcontentsline{toc}{chapter}{Введение: зачем этот фундамент}
\markboth{Введение}{Введение}

Водород, попавший в оболочку твэла из циркониевого сплава, при охлаждении выпадает в виде тонких пластинок
гидрида. Обычно они лежат вдоль окружности трубы и почти не мешают ей держать давление. Но если труба остывает под
растягивающим окружным напряжением, часть пластинок встаёт поперёк --- по радиусу (рис.~\ref{fig:f0_tube}).
Радиальные гидриды режут стенку поперёк и делают металл хрупким. Модель, ради которой написан этот текст,
предсказывает, сколько гидридов переориентируется при заданных водороде, температуре, напряжении и структуре
металла.

\fig{f0_tube}{Переориентация гидридов в сечении трубы (схема). Слева --- до опыта: пластинки гидрида (синие)
лежат вдоль окружности. Справа --- после охлаждения под окружным напряжением $\sigma_\theta$: часть пластинок
встала по радиусу (красные).}

Чтобы такую модель построить или хотя бы понять, нужно одновременно держать в голове несколько областей:
строение металлов и кристаллографию, методы исследования, термодинамику растворов и фазовые диаграммы, диффузию,
теорию упругости, теорию зарождения, механику и пластичность. Каждая из них --- отдельный учебник. Здесь собрано
то, что реально работает в задаче о гидридах, и объяснено так, чтобы было видно, \emph{почему} формулы такие, а
не только \emph{какие} они.

\section*{Как устроен текст}

Текст разбит на четыре части. Первая --- азбука материаловедения: из чего состоит металл и как это видят приборы.
Вторая --- равновесие: какие фазы устойчивы и сколько их. Третья --- как фазы появляются и растут. Четвёртая ---
как металл сопротивляется нагрузке и как его ломают гидриды.

\begin{center}
\small
\begin{tabularx}{\linewidth}{>{\raggedright\arraybackslash}p{4.6cm}X>{\centering\arraybackslash}p{1.4cm}}
\toprule
что есть в модели & что нужно понимать & глава \\
\midrule
\multicolumn{3}{l}{\textit{Часть I. Строение твёрдого тела}} \\
упругие модули, тепловое расширение, водород как частица в решётке & связь атомов, колебания решётки, электронные зоны & \ref{ch:atoms} \\
решётка $\alpha$-Zr, гидрид, плоскость залегания & ячейки, плотные упаковки, индексы Миллера--Браве, пустоты, стереографическая проекция & \ref{ch:lattice} \\
места зарождения, механизм образования гидрида, границы & вакансии, дислокации, системы скольжения ГПУ, частичные дислокации, межфазные границы & \ref{ch:defects} \\
исходные данные и сравнение с опытом & рентген и дифракция, синхротрон, EBSD, металлография и травление, ДСК, анализ водорода & \ref{ch:methods} \\
ориентации зёрен трубы & полюсные и обратные полюсные фигуры, ФРО, факторы Кернса, происхождение текстуры & \ref{ch:texture} \\
\multicolumn{3}{l}{\textit{Часть II. Термодинамика и фазовые равновесия}} \\
пределы растворимости водорода, гистерезис & энергия Гиббса, химический потенциал, растворы, общая касательная & \ref{ch:gibbs}, \ref{ch:solutions} \\
сколько гидрида при данной температуре & правило фаз, правило рычага, эвтектика и эвтектоид, диаграмма Zr--H & \ref{ch:diagrams} \\
\multicolumn{3}{l}{\textit{Часть III. Диффузия и превращения}} \\
перенос водорода к растянутым местам & законы Фика, химический потенциал как движущая сила, эффект Горского & \ref{ch:diffusion} \\
энергия пластинки, поле напряжений, работа нагрузки & собственная деформация, задача Эшелби, метод Фурье & \ref{ch:elastic} \\
когда и где появляется пластинка, как растёт & классическая теория зарождения, неоднородное зарождение, кинетика & \ref{ch:nucleation} \\
\multicolumn{3}{l}{\textit{Часть IV. Механика}} \\
зачем всё это: опасность радиальных гидридов & испытания, механика разрушения, ползучесть и релаксация & \ref{ch:mech} \\
пластическая аккомодация, память, анизотропия трубы & текучесть, факторы Шмида, упрочнение, рекристаллизация & \ref{ch:plasticity} \\
\bottomrule
\end{tabularx}
\end{center}

Главы можно читать по порядку: каждая опирается на предыдущие. Если материаловедение знакомо, часть~I можно
пролистать и вернуться к ней по ссылкам. Если нет --- лучше начать с неё: без языка решёток, индексов и полюсных
фигур дальнейшее читается с трудом. В каждой главе есть блоки четырёх видов.
\begin{itemize}
\item \textbf{\color{circ}Интуиция} --- что происходит «на пальцах», до формул.
\item \textbf{\color{load}Главная формула} --- то, что стоит запомнить или уметь вывести.
\item \textbf{\color{hydr}Для гидридов} --- как общий результат применяется к водороду в цирконии, с числами.
\item \textbf{\color{rad}Осторожно} --- где легко ошибиться.
\end{itemize}
В конце глав --- упражнения; ответы и подсказки собраны в приложении Б. Справочные величины для системы Zr--H
с источниками --- в приложении А. Все рисунки посчитаны по формулам этого текста (скрипты лежат в
\texttt{figs\_src/}); числа на них можно проверить самостоятельно. Схемы, где расчёт заменён рисунком, так и
подписаны.

\section*{Единицы и масштабы}

В задаче о гидридах встречаются энергии «на атом», «на моль» и «на кубический метр»; их полезно уметь переводить
друг в друга с ходу.
\begin{itemize}
\item Тепловая энергия: $k_BT$ на атом или $RT$ на моль. При 300\degC{} ($T = 573$~К)
  $RT \approx 4.8$~кДж/моль, $k_BT \approx 0.049$~эВ. Это мера «тепловых толчков»: барьер в $20RT$ при этой
  температуре преодолевается редко, в $2RT$ --- постоянно.
\item Концентрацию водорода в цирконии обычно дают в \emph{ppm по массе} (wppm, млн$^{-1}$): микрограммы водорода
  на грамм металла. Атомная доля $x_{\mathrm H} \approx C_{\mathrm{ppm}} \cdot 91.22 / (1.008\cdot 10^6)$, то есть
  100~ppm $\approx 0.9$~ат.\,\%.
\item Энергия на единицу объёма: 1~МДж/м$^3$ = 1~МПа. Напряжение --- это и есть плотность энергии. Поэтому
  работа нагрузки 100~МПа на деформации 7\,\% равна 7~МДж/м$^3$, и её можно прямо сравнивать с упругой энергией
  гидрида в тех же единицах (глава~\ref{ch:elastic}).
\item Объём гидрида на моль водорода около $10^{-5}$~м$^3$. Поэтому 1~МДж/м$^3$ гидрида $\approx 10$~Дж на моль водорода.
\end{itemize}
